(1) 一摩尔气块焓为 $h=C_PT$，由可逆热力学第一定律 $Tds=dh-v_mdp=0$，得 $C_PdT=v_mdp$。周围大气静力平衡 $dp/dz=-\rho g=-Mg/v_m$，因此沿机械平衡的抬升路径 $C_PdT/dz=-Mg$。

(2) 由上式得到 $\boxed{dT/dz=-Mg/C_P\equiv-\Gamma_d}$。理想气体摩尔体积 $v_m=RT/p$，故 $C_P dT=(RT/p)dp$，积分为 $\boxed{pT^{-C_P/R}=\mathrm{const}}$，与经典Poisson绝热关系一致。

(3) 将气块由 $z$ 向上扰动 $\delta z>0$，气块绝热温度变化 $-\Gamma_d\delta z$，环境变化 $T_e'\delta z$，两者温差为 $\delta T=(-\Gamma_d-T_e')\delta z$。在近似相等压力下温度高则密度小；因此气块受到恢复性向下浮力要求 $-\Gamma_d-T_e'<0$，即 $\boxed{T_e'>-\Gamma_d}$，也就是环境直减率小于干绝热直减率。

(4) 若全大气干绝热且 $C_P$ 常数，则 $T(z)=T_0-\Gamma_d z$，$p(z)=p_0[T(z)/T_0]^{C_P/R}$。当 $z$ 逼近 $T_0/\Gamma_d$ 时模型给零绝对温度，实际早在此处之前热容、湿度、辐射、成分或非连续平衡条件已发生变化，不能向负温区外推。
